\section{Lucas-carrier arithmetic and redundancy ten}\label{sec:lucas}

This appendix evaluates the modular carrier left by the recursive theorem.
It first separates the additive endpoint orbit from the full degree-nine
carrier, then proves that the full carrier is shallow, and finally combines
that arithmetic with five pointed contractions to obtain redundancy ten.

Let $d=2^s$ with $s\geq2$ in characteristic two.  Lucas' theorem gives the top
nucleus of the degree-$d$ normal rational curve as
\[
 \PP\langle e_1,\ldots,e_{d-1}\rangle,
\]
whose coherent degree-$(d+1)$ lift is
\[
 \cM_{d+1}=\PP\langle e_2,\ldots,e_{d-1}\rangle.
\]
At the distinguished Borel endpoint $e_{d-1}$ the Hankel kernel
contains
\[
       \cU_s=\langle1,t,t^d\rangle.
\]
We use the standard linearized/subspace-polynomial terminology of
\cite{LidlNiederreiter1996,BenSassonEtAl2016}; the carrier and
root-cover calculations below are specific to the present Hankel
system.

\begin{proposition}[Lucas overlap kernel]\label{prop:lucas-kernel}
For $d=2^s$, the lower nucleus and its coherent lift are exactly
\[
\PP\langle e_1,\ldots,e_{d-1}\rangle,\qquad
\PP\langle e_2,\ldots,e_{d-1}\rangle.
\]
At $e_{d-1}$, $\cU_s$ is endpoint-specific; the intersection common to
the full carrier is only $\langle1,t^d\rangle$.
\end{proposition}

\begin{proof}
Lucas' theorem gives
$\binom di\equiv0\pmod2$ for $0<i<d$.  Applying this support condition
to the two consecutive contraction rows removes the two endpoint
coordinates and gives the displayed lift.  For a basis point $e_i$,
the Hankel equations kill the adjacent coefficients $b_{i-1},b_i$.
At $i=d-1$, the monomials $1,t,t^d$ survive.  Varying $i$ over the
carrier leaves only $1,t^d$ in the intersection.
\end{proof}

\begin{proposition}[Linearized root cover]\label{prop:linearized}
On the generic locus $BC\ne0$, the normalized ordered-root cover of
$A+Bt+Ct^d$ has geometric
monodromy $\operatorname{AGL}_1(\F_d)$, minimal constant field
$\F_d$, and based deck group $C_{d-1}$.  Coefficientwise Frobenius acts
on the based group by multiplication by two, with exact order $s$.
The net contains a split squarefree member over $\F_{2^m}$ if and only
if $s\mid m$.
\end{proposition}

\begin{figure}[H]
\centering
\[
\begin{array}{ccccc}
K=\overline{\F}_2(a,b)
&\subset&
K_1=K(\beta),\quad \beta^{d-1}=b
&\subset&
K_1(y),\quad y^d+y=a/\beta^d\\[3pt]
&&\underbrace{\F_d^\times}_{\text{scalings},\ d-1}
&&\underbrace{(\F_d,+)}_{\text{translations},\ d}\\[5pt]
&&\multicolumn{3}{c}{
(\F_d,+)\rtimes\F_d^\times=\operatorname{AGL}_1(\F_d)}
\end{array}
\]
\caption{The linearized ordered-root cover.  The Kummer scaling layer
and connected additive layer have coprime roles; their explicit
semidirect action accounts for the full geometric monodromy.}
\label{fig:lucas-monodromy}
\end{figure}

\begin{proof}
Work first over the geometric coefficient field
\[
 k=\overline{\F}_2,\qquad K=k(A/C,B/C),
\]
on the open set $BC\ne0$.  Put $a=A/C$ and $b=B/C$.  The Kummer
extension
\[
 K_1=K(\beta),\qquad \beta^{d-1}=b,
\]
has degree $d-1$: the valuation of $b$ at its simple zero is one.
Since $k$ contains $\F_d$, this is the scaling part of the geometric
root cover, not an extension of constants.  Over $K_1$ choose
\[
 y^d+y=a/\beta^d.
\]
The roots are $\{\beta(y+c):c\in\F_d\}$.  The right side has a simple
pole at the pole of $a$.  If the $\F_d$-torsor were disconnected, its
geometric translation group would lie in a proper additive subgroup
of $\F_d$.  A nonzero $\F_2$-linear functional annihilating that
subgroup would then give an Artin--Schreier equation
$w^2+w=\gamma a/\beta^d$ with $\gamma\in\F_d^\times$.  This is
impossible: the right side has a simple pole, whereas every pole of
$w^2+w$ has even order.  Thus $K_1(y)/K_1$ is connected of degree
$d$.
Its translations $y\mapsto y+c$ give the normal subgroup
$(\F_d,+)$.  The $d-1$ Kummer conjugates act by
\[
 (\beta,y)\longmapsto(\zeta\beta,\zeta^{-1}y),
 \qquad \zeta\in\F_d^\times,
\]
and conjugate translations by multiplication by $\zeta$.  We have
thus exhibited $d(d-1)$ automorphisms of an extension of degree
$d(d-1)$.  Consequently the full geometric group, rather than a
proper subgroup, is
\[
 (\F_d,+)\rtimes\F_d^\times
   =\operatorname{AGL}_1(\F_d).
\]

Return now to the arithmetic coefficient field
$\F_2(A/C,B/C)$.  Ratios of nonzero root differences recover every
element of $\F_d$, so any splitting field contains $\F_d$ in its
constant field.  The preceding regular construction over $\F_d$
introduces no larger constants, proving minimality.  Coefficient
Frobenius acts on $\F_d^\times$ by $c\mapsto c^2$, and hence has exact
order $s$ on the based component group $C_{d-1}$.  Boundary members
with $C=0$ have degree at
most one, while those with $B=0$ have zero derivative; neither is split
squarefree of degree $d$.  Thus a split member forces
$\F_d\subset\F_{2^m}$, equivalently $s\mid m$.  Conversely, $t^d+t$
splits simply under that inclusion.  The equality
$\operatorname{ord}_{2^s-1}(2)=s$ follows from
\[
\gcd(2^r-1,2^s-1)=2^{\gcd(r,s)}-1.
\]
\end{proof}

At $s=2$ this recovers the odd/even extension law at redundancies six
and seven.  The first fresh carrier is $s=3$, in syndrome degree nine.
The $\F_8$-linear section has order-three component Frobenius, but its
four quotient directions change the full answer.

\begin{proposition}[The $e_7$ orbit]\label{prop:e7-orbit}
The stabilizer of $e_7\in\PP(\Gamma^9E)$ is the upper Borel, so its
$\PGL_2$-orbit is $\PP^1$.  At the distinguished point,
\[
W_{e_7}=\langle1,t,t^2,t^3,t^4,t^5,t^8\rangle.
\]
\end{proposition}

\begin{proof}
Lucas parity and extraction of the $x^2y^7$ coefficient give
\[
g e_7=(ad+bc)^2
\bigl(b^5e_2+b^4d\,e_3+bd^4e_6+d^5e_7\bigr).
\]
This is proportional to $e_7$ exactly when $b=0$.  The apolar action is
contragredient, so it transports both $W_f$ and split root divisors.
The displayed kernel follows from the two Hankel equations.
\end{proof}

\begin{proposition}[Framed Artin--Schreier quotient]
\label{prop:e7-as}
Order roots $0,1,a,b,c,d,u,v$ and put
\[
S=a+b+c+d,\quad E=ab+ac+ad+bc+bd+cd,\quad
h=1+S,\quad p=1+E+S+S^2.
\]
On the open set where the six fixed roots are distinct,
$h,p,1+h+p$ are nonzero, and
$x^2+hx+p\ne0$ for $x=a,b,c,d$, the remaining pair is the
geometrically integral Artin--Schreier cover
\[
       y^2+y=p/h^2.
\]
Its rational lifting law is
$\Tr_{\F_{2^m}/\F_2}(p/h^2)=0$.
\end{proposition}

\begin{proof}
The first elementary-symmetric equation gives
$u+v=1+S=h$.  In the second equation the root $1$ contributes
$S+u+v=1$, so
\[
 0=E+S(u+v)+uv+S+u+v
   =E+Sh+uv+1,
\]
and hence $uv=p$.  Thus $u,v$ are the roots of
$z^2+hz+p$.  The divisor $h=0$ is exactly their collision divisor.
On $h\ne0$, putting $y=u/h$ gives the displayed
Artin--Schreier equation.

It remains to prove that its class is geometric, rather than a
constant-field artifact.  Write $B=b+c+d$, so $h=a+1+B$.  Along the
prime divisor $h=0$, the leading coefficient of the order-two pole of
$p/h^2$ is
\[
D=1+bc+bd+cd+B+B^2,
\]
in the residue field $\overline{\F}_2(b,c,d)$.  Its derivative
$\partial D/\partial b=1+c+d$ is a nonzero rational function, so $D$
is not a square there.  If $p/h^2=g^2+g$, the order-two pole would
force $g$ to have a simple pole whose principal coefficient squares
to $D$, a contradiction.  The class therefore stays nonzero after
algebraic extension of the constants, proving geometric integrality
and geometric monodromy $C_2$.

The remaining open conditions have exact collision meanings:
$p\ne0$ excludes $u,v=0$, $1+h+p\ne0$ excludes $u,v=1$, and
$x^2+hx+p\ne0$ excludes $u,v=x$ for
$x\in\{a,b,c,d\}$.  Over $\F_{2^m}$ the standard
Artin--Schreier criterion gives a rational pair precisely when
$\Tr_{\F_{2^m}/\F_2}(p/h^2)=0$.
\end{proof}

\begin{proposition}[Additive frame cover]\label{prop:e7-additive}
The family
\[
t^8+A_4t^4+A_2t^2+A_1t+A_0,\qquad A_1\ne0,
\]
has a connected affine-frame cover with geometric monodromy
$\operatorname{AGL}_3(\F_2)$.
\end{proposition}

\begin{proof}
For an $\F_2$-space $U$ and $w\notin U$,
\[
L_{U+\langle w\rangle}=L_U^2+L_U(w)L_U.
\]
Induction gives exponents $8,4,2,1$ and
$A_1=\prod_{v\ne0}v\ne0$.  An affine frame is an origin plus three
independent differences.

Let $\mathcal F$ be the open affine frame space of tuples
$(a;v_1,v_2,v_3)$ with independent $v_i$.  It is geometrically
integral.  Sending a frame to the monic polynomial with root set
\[
 a+\langle v_1,v_2,v_3\rangle_{\F_2}
\]
lands in the displayed coefficient family.  Conversely, the eight
geometric roots of a generic member form an affine three-space: the
difference set is the kernel of its additive part, and any root is an
origin.  Its ordered affine frames are therefore exactly one free
orbit of
\[
 \F_2^3\rtimes\operatorname{GL}_3(\F_2).
\]
Thus the generic fiber has exactly
$8|\operatorname{GL}_3(\F_2)|=1344$ points, every two of them differ
by one affine transformation, and the invariant function field is the
coefficient field.  The normalization is connected and has exact
geometric deck group $\operatorname{AGL}_3(\F_2)$.
\end{proof}

\begin{remark}[Logical boundary]
The generic Artin--Schreier cover describes the full framed quotient,
and Proposition~\ref{prop:e7-additive} identifies a proper additive
subcover.  The shallowness theorem below does not require a rational
point on the generic cover: its witnesses are constructed directly
from rational three-dimensional subspaces.
\end{remark}

\begin{theorem}[Shallowness of the degree-nine $e_7$ orbit]
\label{thm:e7}
For every $m\geq3$, the entire $\PGL_2$-orbit of $e_7$ is shallow over
$\F_{2^m}$.  The number of monic additive-affine witnesses is
\[
\frac{q(q-1)(q-2)(q-4)}{1344}.
\]
\end{theorem}

\begin{proof}
Every three-dimensional $\F_2$-subspace $V\subset\F_{2^m}$ gives
\[
L_V(t)=\prod_{v\in V}(t-v)
      =t^8+A_4t^4+A_2t^2+A_1t.
\]
The $t^7,t^6$ coefficients vanish, so $L_V\in W_{e_7}$, and
$A_1\ne0$ makes its roots simple.  Every admissible binary
field in the statement contains such a three-space.
Proposition~\ref{prop:e7-orbit} transports the witness around the
orbit.

There are
\[
\binom m3_2=\frac{(q-1)(q-2)(q-4)}{168}
\]
three-spaces.  Each has $q/8$ affine cosets, and
$L_{a+V}(t)=L_V(t)+L_V(a)$.  A split polynomial recovers its root
coset, whose difference set recovers $V$; hence distinct cosets or
three-spaces cannot be counted twice.  Multiplying the two factors
gives the displayed formula.
\end{proof}

Thus the order-three law of Proposition~\ref{prop:linearized} belongs
to a proper $\F_8$-linear section, not to deepness of the full
$e_7$ kernel.  We now treat the complete carrier.

\subsection{The complete degree-nine carrier}

Write
\[
 \cM_9=\PP\langle e_2,e_3,e_4,e_5,e_6,e_7\rangle,\qquad
 U=\langle e_2,e_3,e_6,e_7\rangle .
\]
The $\PGL_2$-module filtration gives
\[
                    \cM_9/U\simeq\det^4\otimes E.
 \autoeqtag{eq:lucas-m9-quotient}
\]

\paragraph{Proof route}
The proof of \(\cM_9(\F_{2^m})\) shallowness has four distinct steps.
\begin{enumerate}[label=\textup{(\arabic*)}]
\item On the invariant block \(\PP(U)\), reciprocal subspace polynomials
      settle the rank-one strata.  The remaining rank-two twists reduce to
      the Artin--Schreier final-pair cover in
      Proposition~\ref{prop:m9-invariant-block}.
\item Off \(\PP(U)\), projective transport fixes the nonzero quotient
      coordinate.  The coefficient map in \eqref{eq:binary-hankel-rank} has
      rank three or four; in
      either case the nonsquare double-pole calculation makes the
      Artin--Schreier cover geometrically integral.
\item Fixing five roots leaves one moving base root.  The good-base
      conditions have bounded degree, and the normalized cover has genus at
      most one with deletion degree \(48\).  This gives every complement
      point for \(q\geq64\).
\item The finite boundary is explicit.  Full-carrier quotient certificates
      close \(q=16,32\).  At \(q=64\), only the invariant block is
      certificate-backed; the complement is covered by the preceding
      mathematical point-count argument.  For \(q\geq128\), both parts are
      mathematical.
\end{enumerate}
Thus the rank and selector calculations carry the uniform theorem, while the
certificates close exactly the stated finite residue.

\begin{proposition}[Invariant-block arithmetic]\label{prop:m9-invariant-block}
Every rational point of $\PP(U)$ is shallow over every
$\F_{2^m}$ with $m\geq4$.
\end{proposition}

\begin{proof}
Identify \(U\) with \(2\times2\) matrices by
\(e_2,e_3,e_6,e_7\leftrightarrow E_{00},E_{01},E_{10},E_{11}\).
A rank-one matrix is \(u^{(4)}v^{\mathsf T}\).  The diagonal
\([u]=[v]\) is the Frobenius-graph orbit and is handled by the subspace
polynomial of a three-dimensional \(\F_2\)-space.  The off-graph orbit has
representative \(e_3\), whose kernel equations are \(f_2=f_3=0\).  If
\[
 L_V(X)=X^8+A_4X^4+A_2X^2+A_1X,\qquad A_1\ne0,
\]
is the subspace polynomial of a three-space \(V\subset\F_{2^m}\), then
\[
 A_1^{-1}(t+A_4t^5+A_2t^7+A_1t^8)
\]
has the eight distinct roots
\(\{0\}\cup\{v^{-1}:0\ne v\in V\}\) and satisfies those equations.
Projective transport closes both rank-one strata.

For rank two write the rational syndrome as
\(z=(a,b,c,d)\), with \(ad+bc\ne0\).  This is the rational
\(A_5\)-twist family of \(e_3+e_6\), but the construction below is defined
over \(z\) itself and requires no chosen cocycle.  Choose six roots with
polynomial $h=\sum_{i=0}^6h_it^i$, and put
\[
\begin{array}{ll}
A_0=bh_0+ch_3+dh_4,&A_1=ah_0+bh_1+ch_4+dh_5,\\
A_2=ah_1+bh_2+ch_5+d,&A_3=ah_2+bh_3+c.
\end{array}
\]
and set
\[
 \Delta=A_1A_3+A_2^2,\quad Q=A_0A_3+A_1A_2,\quad
 N=A_1^2+A_0A_2.
\]
The final pair has sum $Q/\Delta$, product $N/\Delta$, and rational
distinct roots precisely when
\[
                         y^2+y=N\Delta/Q^2
 \autoeqtag{eq:artin-schreier-cover}
\]
has a rational point away from the root-collision divisors.  These equations
descend the cover uniformly across all twists.  Fix five roots and let \(x\)
be the sixth.  Then \(Q\) has degree at most two in \(x\), \(N\Delta\) has
degree at most four, and \(Q'\) is a nonzero polynomial in the five-root
parameters; the good-base open below imposes \(Q'\ne0\).  At a root of \(Q\),
Artin--Schreier reduction leaves a simple pole exactly when
\[
 R=((N\Delta)')^2+(Q')^2N\Delta\ne0.
\]
At the alternating representative \(e_3+e_6\), the four \(A_i\) are
independent; on \(Q=0\) and \(A_3\ne0\) one has
\[
 A_0=A_1A_2/A_3,\qquad N\Delta=\Delta^2A_1/A_3.
\]
Thus the double-pole leading coefficient is \(A_1/A_3\), which is not a
square in the residue function field.  Every geometric rank-two point is
projectively conjugate to that representative, so the induced root-coordinate
change is birational.
Ordering the roots is finite separable; it cannot make the coefficient a
square, because adjoining its square root in characteristic two would be
purely inseparable.  Thus the same nonsquare test holds on every rational
twist and says \(Q\nmid R\).  The cover is geometrically integral and has at
most two reduced simple poles, so genus at most one.  A nonzero good-base
selector has total degree at most \(102\) in the five ordered roots.  For a
good base the forbidden
\(x\)-values have degrees \(5,2,2,10,4\), totaling \(23\), and therefore
delete at most \(2\cdot23+2=48\) points of the smooth double cover.  The
zero bound gives a good base for \(q\ge128\), while Hasse--Weil then supplies
a surviving point.
The complete rank-two invariant-block quotient certificate at $q=64$ and the
$q=16,32$ quotient certificates cover the remaining admissible fields.
Their checkers rebuild
the orbit partitions, twists, root polynomials, and both Hankel equations;
the generator and separately written replay are the
\texttt{rank-two-artin-schreier-avoidance} and
\texttt{higher-lucas-modular-carriers} artifacts under
\texttt{supplement/evidence/lucas-m9/}.  Section~\ref{sec:verification}
records their exact trust boundary.
\end{proof}

\begin{proposition}[Exact final-pair criterion]\label{prop:final-pair}
Let $z=(z_2,\ldots,z_7)\in\cM_9(\F_{2^m})$.  Choose six distinct
finite roots with monic polynomial $h=\sum_{i=0}^6h_it^i$, and define
\[
 A_j=\sum_{i=2}^7z_i h_{i+j-3}\quad(0\leq j\leq3),
\]
with coefficients outside $0,\ldots,6$ equal to zero.  If the final
pair has sum $s$ and product $p$, the two Hankel equations are exactly
\[
 A_0+sA_1+pA_2=0,\qquad A_1+sA_2+pA_3=0.
 \autoeqtag{eq:final-pair-linear}
\]
On
\[
 \Delta=A_1A_3+A_2^2\ne0,\qquad
 Q=A_0A_3+A_1A_2\ne0,
\]
their unique solution is
\[
 s=Q/\Delta,\qquad p=(A_1^2+A_0A_2)/\Delta,
 \autoeqtag{eq:final-pair-sp}
\]
and the pair is rational and distinct exactly when
\[
 \Tr_{\F_{2^m}/\F_2}\!
 \left(\frac{(A_1^2+A_0A_2)\Delta}{Q^2}\right)=0.
 \autoeqtag{eq:final-pair-trace}
\]
Collision with an old root $r$ is
\[
 r^2\Delta+rQ+(A_1^2+A_0A_2)=0.
 \autoeqtag{eq:final-pair-quadratic}
\]
The rank-deficient cases are exactly the corresponding rank-one or
rank-zero alternatives in \eqref{eq:final-pair-linear}.
\end{proposition}

\begin{proof}
Substitute $f=h(t)(t^2+st+p)$ in the two consecutive Hankel equations.
This gives \eqref{eq:final-pair-linear}; Cramer's rule gives
\eqref{eq:final-pair-sp}.  In characteristic two the
quadratic is split with distinct roots precisely when $s\ne0$ and
$\Tr(p/s^2)=0$, which is \eqref{eq:final-pair-trace}.  Evaluation of the
quadratic at $r$ gives \eqref{eq:final-pair-quadratic}.  Conversely, every
split divisor of nonfull support can be
moved off infinity and partitioned into six roots and a final pair.
The full-support divisor is checked directly in homogeneous coordinates.
Thus the criterion loses neither infinity nor the degenerate linear-system
cases.
\end{proof}

The next three lemmas carry the all-field part of the theorem below: the
first fixes the rank of the coefficient map off the invariant block, the
second makes the final-pair cover geometrically integral, and the third
selects a rational base of five roots with an effective degree bound.

\begin{lemma}[Rank of the coefficient map off the invariant block]
\label{lem:m9-rank}
Let $z=(z_2,\ldots,z_7)$ and let $h=(h_0,\ldots,h_5)^{\mathsf T}$ collect the
free coefficients of the monic sextic of Proposition~\ref{prop:final-pair}.
Then
\[
 (A_0,A_1,A_2,A_3)^{\mathsf T}=M_zh+(0,0,z_7,z_6)^{\mathsf T},
\quad
M_z=\begin{pmatrix}
 z_3&z_4&z_5&z_6&z_7&0\\
 z_2&z_3&z_4&z_5&z_6&z_7\\
 0&z_2&z_3&z_4&z_5&z_6\\
 0&0&z_2&z_3&z_4&z_5
\end{pmatrix},
\autoeqtag{eq:binary-hankel-rank}
\]
and $\operatorname{rank}M_z\geq3$ whenever $z_4\neq0$ or $z_5\neq0$.  In that
case exactly one of the following holds.
\begin{enumerate}[label=\textup{(\roman*)}]
\item $\operatorname{rank}M_z=4$, and $h\mapsto(A_0,A_1,A_2,A_3)$ is onto
      $\mathbb A^4$.
\item $\operatorname{rank}M_z=3$, the left kernel of $M_z$ is a line
      $\langle\ell\rangle$, and the image of $h\mapsto(A_0,A_1,A_2,A_3)$ is
      the affine hyperplane
      \[
       \ell_0A_0+\ell_1A_1+\ell_2A_2+\ell_3A_3=\kappa,
       \qquad\kappa=\ell_2z_7+\ell_3z_6 .
       \autoeqtag{eq:m9-affine-relation}
      \]
\end{enumerate}
\end{lemma}

\begin{proof}
Reading $A_j=\sum_{i=2}^7z_ih_{i+j-3}$ row by row, with $h_6=1$ and $h_k=0$
for $k\notin\{0,\ldots,6\}$, gives the displayed affine form; the constant
vector collects the two $h_6$ terms, which occur in $A_2$ and $A_3$ only.
A row vector $\ell$ satisfies $\ell M_z=0$ exactly when
\[
\begin{aligned}
 \ell_0z_3+\ell_1z_2&=0,&
 \ell_0z_4+\ell_1z_3+\ell_2z_2&=0,\\
 \ell_0z_5+\ell_1z_4+\ell_2z_3+\ell_3z_2&=0,&
 \ell_0z_6+\ell_1z_5+\ell_2z_4+\ell_3z_3&=0,\\
 \ell_0z_7+\ell_1z_6+\ell_2z_5+\ell_3z_4&=0,&
 \ell_1z_7+\ell_2z_6+\ell_3z_5&=0,
\end{aligned}
\autoeqtag{eq:m9-column-relations}
\]
one equation per column, read left to right.

Suppose $\operatorname{rank}M_z\leq2$, so the left kernel $K$ has dimension at
least two.  Either $K$ contains a nonzero vector with $\ell_0=\ell_1=0$, or
$\ell\mapsto(\ell_0,\ell_1)$ is injective on $K$ and hence, $\dim K\geq2$,
also onto.

In the first case take $\ell=(0,0,\ell_2,\ell_3)\neq0$.  If $\ell_2\neq0$,
normalize $\ell_2=1$: the second, third, fourth and fifth equations of
\eqref{eq:m9-column-relations} read $z_2=0$, $z_3=\ell_3z_2$,
$z_4=\ell_3z_3$ and $z_5=\ell_3z_4$, so $z_2=z_3=z_4=z_5=0$ in turn.  If
$\ell_2=0$ then $\ell_3\neq0$, and the third through sixth equations read
$\ell_3z_2=\ell_3z_3=\ell_3z_4=\ell_3z_5=0$, with the same conclusion.

In the second case $K$ contains $(1,0,\ell_2,\ell_3)$ and
$(0,1,\ell_2',\ell_3')$.  The first equation applied to the second vector
gives $z_2=0$, and applied to the first gives $z_3=0$; the second and third
equations applied to the first vector then give $z_4=\ell_2z_2=0$ and
$z_5=\ell_2z_3+\ell_3z_2=0$.

Either case contradicts $z_4\neq0$ or $z_5\neq0$, so
$\operatorname{rank}M_z\in\{3,4\}$.  At rank four the linear map $h\mapsto
M_zh$ is onto, hence so is its translate.  At rank three the left kernel is a
line $\langle\ell\rangle$, the image of $h\mapsto M_zh$ is the linear
hyperplane $\sum_i\ell_iA_i=0$, and translating by $(0,0,z_7,z_6)$ moves the
right-hand side to $\kappa=\ell_2z_7+\ell_3z_6$.
\end{proof}

\begin{lemma}[Geometric nontriviality of the final-pair Artin--Schreier
cover]\label{lem:m9-nonsquare}
Let $k$ be algebraically closed of characteristic two, let $z\in\cM_9(k)$
satisfy $z_4\neq0$ or $z_5\neq0$, and let $L\subseteq\mathbb A^4$ be the
image of the coefficient map of Lemma~\ref{lem:m9-rank}.  Then the reduced
divisor $Q=0$ in $L$ has an irreducible component $\mathfrak q$ on which
$A_3$ and $\Delta$ do not vanish identically and on whose function field
$A_1/A_3$ is not a square.  Ordering the six roots leaves that conclusion
intact, so the double pole of $N\Delta/Q^2$ along $\mathfrak q$ does not
reduce and \eqref{eq:artin-schreier-cover} is geometrically integral.
\end{lemma}

\begin{proof}
On $Q=0$ with $A_3\neq0$ one has $A_0=A_1A_2/A_3$, hence
\[
 N=A_1^2+\frac{A_1A_2^2}{A_3}=\frac{A_1(A_1A_3+A_2^2)}{A_3}
  =\frac{A_1\Delta}{A_3},
 \qquad
 N\Delta=\frac{A_1}{A_3}\,\Delta^2 .
 \autoeqtag{eq:m9-double-pole}
\]
Since $u^2\equiv u$ modulo $\wp(u)=u^2+u$, the double pole of $N\Delta/Q^2$
along a component of $Q=0$ reduces to a simple pole exactly when $A_1/A_3$ is
a square in the function field of that component.  That square class is what
the two rank cases test.

\emph{Rank four.}  Here $L=\mathbb A^4$ and the coordinates $A_i$ are
independent.  The partial derivatives of $Q=A_0A_3+A_1A_2$ are
$A_3,A_2,A_1,A_0$, whose only common zero is the origin, so the projective
quadric is smooth and $X=V(Q)$ is irreducible; in particular $A_3$ and
$\Delta$ do not vanish on it identically.  Inside $X$ lie the planes
$\Pi_1=V(A_0,A_1)$, $\Pi_{13}=V(A_1,A_3)$ and $\Pi_2=V(A_2,A_3)$, and
$V(A_1)\cap X=\Pi_1\cup\Pi_{13}$, $V(A_3)\cap X=\Pi_{13}\cup\Pi_2$.  At the
generic point of $\Pi_1$ the functions $A_2,A_3$ are units and
$A_0=A_1A_2/A_3$, so the local ring of $X$ there is a discrete valuation ring
with uniformizer $A_1$; symmetrically $A_3$ is a uniformizer at $\Pi_2$,
where $A_0,A_1$ are units and $A_2=A_0A_3/A_1$; and at $\Pi_{13}$ the two
relations $A_3=A_1A_2/A_0$ and $A_1=A_0A_3/A_2$ show that $A_1$ and $A_3$
generate the same maximal ideal.  Hence
\[
 \operatorname{div}\left(\frac{A_1}{A_3}\right)=[\Pi_1]-[\Pi_2],
\]
with both multiplicities equal to one.  A square has even multiplicity along
every prime divisor, so $A_1/A_3$ is not a square in $k(X)$, and
$\mathfrak q=X$ serves.

\emph{Rank three, principal case.}  Let \eqref{eq:m9-affine-relation} be the
relation and assume first $(\ell_0,\ell_1)\neq(0,0)$.  We claim some
component $\mathfrak q$ of $Q=0$ in $L$ carries $A_2$ and $A_3$ as
algebraically independent functions.  Otherwise every component lies in an
affine plane $\alpha A_2+\beta A_3=\gamma$ with $(\alpha,\beta)\neq(0,0)$.
Such a plane is not all of $L$, because $(0,0,\alpha,\beta)$ is not
proportional to $\ell$; and a component of the quadric surface $Q|_L$ is
either that whole surface, which has degree two and so lies in no plane, or a
plane.  Hence every component is one of the planes
$\alpha A_2+\beta A_3=\gamma$, and $Q|_L$, being of degree two, is a scalar
multiple of a product of two of those linear forms, hence a polynomial in
$A_2,A_3$ alone.  That is impossible: if $\ell_0\neq0$,
eliminating $A_0$ leaves
$Q|_L=A_3(\kappa+\ell_1A_1+\ell_2A_2+\ell_3A_3)/\ell_0+A_1A_2$, whose
coefficient of $A_1$ is $\ell_1A_3/\ell_0+A_2\neq0$; and if $\ell_0=0$,
$\ell_1\neq0$, eliminating $A_1$ leaves
$Q|_L=A_0A_3+A_2(\kappa+\ell_2A_2+\ell_3A_3)/\ell_1$, which contains
$A_0A_3$.  Fix such a component $\mathfrak q$.  On it $A_3\not\equiv0$ and
$\ell_0A_2+\ell_1A_3\not\equiv0$, since either vanishing is a linear
dependence between $A_2$ and $A_3$.

Substituting $A_0=A_1A_2/A_3$ into \eqref{eq:m9-affine-relation} and clearing
$A_3$ gives $A_1(\ell_0A_2+\ell_1A_3)=A_3(\kappa+\ell_2A_2+\ell_3A_3)$, that
is
\[
 \frac{A_1}{A_3}=\frac{\kappa+\ell_2A_2+\ell_3A_3}
                      {\ell_0A_2+\ell_1A_3}
 \autoeqtag{eq:m9-rank-three-slice}
\]
on $\mathfrak q$.  Together with $A_0=A_1A_2/A_3$ and
\eqref{eq:m9-affine-relation} this writes every coordinate as a rational
function of $A_2,A_3$, so $k(\mathfrak q)=k(A_2,A_3)$.  In characteristic two
a rational function of
two independent variables is a square exactly when both of its partial
derivatives vanish.  Writing $P$ and $D$ for the numerator and denominator of
\eqref{eq:m9-rank-three-slice},
\[
 \frac{\partial}{\partial A_2}\frac PD=\frac{\ell_2D+\ell_0P}{D^2}
 =\frac{(\ell_1\ell_2+\ell_0\ell_3)A_3+\ell_0\kappa}{D^2},
 \qquad
 \frac{\partial}{\partial A_3}\frac PD
 =\frac{(\ell_1\ell_2+\ell_0\ell_3)A_2+\ell_1\kappa}{D^2},
\]
so $A_1/A_3$ is a square on $\mathfrak q$ only if
\[
 \ell_0\kappa=\ell_1\kappa=0,\qquad \ell_1\ell_2=\ell_0\ell_3 .
 \autoeqtag{eq:m9-square-obstruction}
\]

Those equalities force $z_4=z_5=0$.  If $\kappa\neq0$ then
$\ell_0=\ell_1=0$, and the first case of the proof of
Lemma~\ref{lem:m9-rank} applies verbatim to $\ell$.  If $\kappa=0$ and
$\ell_0=0$ then $\ell_1\ell_2=0$; for $\ell_1=0$ the same first case applies,
while for $\ell_2=0$ and $\ell_1\neq0$ the first four equations of
\eqref{eq:m9-column-relations} read $\ell_1z_2=0$, $\ell_1z_3=0$,
$\ell_1z_4+\ell_3z_2=0$ and $\ell_1z_5+\ell_3z_3=0$, giving
$z_2=z_3=z_4=z_5=0$.  Finally let $\ell_0\neq0$; then $\kappa=0$, and scaling
$\ell_0=1$ and writing $\lambda=\ell_1$, $\mu=\ell_2$, $\ell_3=\lambda\mu$,
the first three equations of \eqref{eq:m9-column-relations} give
\[
 z_3=\lambda z_2,\qquad z_4=(\lambda^2+\mu)z_2,\qquad
 z_5=\lambda(\lambda^2+\mu)z_2,
 \autoeqtag{eq:m9-borel-elimination}
\]
the fourth gives $z_6=(\lambda^4+\lambda^2\mu+\mu^2)z_2$, and the fifth
substituted into the sixth gives
$(\lambda^2+\mu)z_6+\lambda^2\mu z_4=(\lambda^2+\mu)^3z_2=0$.  If $z_2=0$
then \eqref{eq:m9-borel-elimination} and the remaining equations force $z=0$;
otherwise $\lambda^2+\mu=0$ and \eqref{eq:m9-borel-elimination} gives
$z_4=z_5=0$.
Both outcomes are excluded, so $A_1/A_3$ is not a square on $\mathfrak q$.

\emph{Rank three, relation in $A_2,A_3$ only.}  Suppose $\ell_0=\ell_1=0$, so
the relation is $\ell_2A_2+\ell_3A_3=\kappa$ and $A_0,A_1$ are unconstrained
on $L$.  If $\ell_2\neq0$, use $(A_0,A_1,A_3)$ as coordinates on $L$ and
substitute $A_2=(\kappa+\ell_3A_3)/\ell_2$:
\[
 Q|_L=A_3A_0+A_1(\kappa+\ell_3A_3)/\ell_2 .
\]
This is of degree one in $A_0$, so it is irreducible unless $A_3$ divides its
constant term, which happens only for $\kappa=0$, where
$Q|_L=A_3(A_0+\ell_3A_1/\ell_2)$.  In the irreducible case the unique
component carries $A_1$ and $A_3$ as independent coordinates; in the split
case the component $A_0=\ell_3A_1/\ell_2$ does.  On such a component
$\partial(A_1/A_3)/\partial A_1=1/A_3\neq0$, so $A_1/A_3$ is not a square.
If instead $\ell_2=0$, then $\ell_3\neq0$ and $A_3=\kappa/\ell_3$ is constant
on $L$.  For $\kappa=0$ the function $A_3$ vanishes identically, so the last
row of $M_z$ vanishes and $z_2=z_3=z_4=z_5=0$, which is excluded; for
$\kappa\neq0$ the surface $Q|_L=\kappa A_0/\ell_3+A_1A_2$ is irreducible with
$A_1,A_2$ independent on it, and $A_1/A_3$ is a nonzero constant multiple of
$A_1$, again not a square.

\emph{The component avoids $\Delta=0$.}  On $A_3\neq0$ the locus $Q=\Delta=0$
is $\{(u^3c,u^2c,uc,c)\}$, the two-dimensional irreducible cone over the
twisted cubic, and on $A_3=0$ it is the plane $A_2=A_3=0$.  The chosen
$\mathfrak q$ has $A_3\not\equiv0$, so if $\mathfrak q$ were contained in
$\Delta=0$ it would equal that cone, forcing
$\ell_0u^3c+\ell_1u^2c+\ell_2uc+\ell_3c=\kappa$ for all $u$ and $c$, hence
$\ell=0$ and $\kappa=0$.  In the rank-four case $\mathfrak q$ has dimension
three and the cone has dimension two.  Either way
$\Delta\not\equiv0$ on $\mathfrak q$, so \eqref{eq:m9-double-pole} really
computes the double-pole coefficient there.

\emph{Passage to ordered roots.}  The coefficient map $h\mapsto A$ is onto
$L$ with linear fibers, so the preimage of $\mathfrak q$ is irreducible and
its function field is a purely transcendental extension of
$k(\mathfrak q)$, which acquires no new square roots of elements of
$k(\mathfrak q)$.  Ordering the six roots is a finite separable base change,
and adjoining a square root in characteristic two is purely inseparable, so
it cannot make $A_1/A_3$ a square either.  Hence the double pole survives and
\eqref{eq:artin-schreier-cover} is geometrically integral.
\end{proof}

\begin{lemma}[Good-base selector and degree bound]\label{lem:m9-selector}
Let $z$ satisfy $z_4\neq0$ or $z_5\neq0$.  Fix five roots $r_1,\ldots,r_5$
with $g(t)=\prod_{i=1}^5(t+r_i)=\sum_{i=0}^5g_it^i$, let $x$ be the sixth
root, and put
\[
 B_j=\sum_{i=2}^7z_ig_{i+j-4}\quad(0\leq j\leq4),
 \qquad\text{so that}\qquad A_j=B_j+xB_{j+1}.
 \autoeqtag{eq:m9-five-root-split}
\]
Then $Q=q_2x^2+q_1x+q_0$ with
\[
 q_2=[x^2]Q=B_1B_4+B_2B_3,\qquad
 q_1=Q'=B_0B_4+B_2^2,\qquad
 q_0=B_0B_3+B_1B_2,
 \autoeqtag{eq:m9-Q-coefficients}
\]
neither $q_2$ nor $q_1$ vanishes identically in $g_0,\ldots,g_4$, and
$N\Delta$ and $R=((N\Delta)')^2+(Q')^2N\Delta$ have degree at most four in
$x$.  Write $R=\sum_{i=0}^4R_ix^i$ and let $\rho=\rho_1x+\rho_0$ be the
pseudo-remainder of $R$ by $Q$, defined by $q_2^3R=SQ+\rho$; explicitly
\[
\begin{aligned}
 \rho_1&=q_2^3R_1+q_2^2q_1R_2+(q_2^2q_0+q_2q_1^2)R_3+q_1^3R_4,\\
 \rho_0&=q_2^3R_0+q_2^2q_0R_2+q_2q_1q_0R_3+(q_2q_0^2+q_1^2q_0)R_4 .
\end{aligned}
 \autoeqtag{eq:m9-pseudo-remainder}
\]
Some $\rho_i$ is a nonzero polynomial, and the product of that $\rho_i$ with
$q_2$, $q_1$ and the Vandermonde $\prod_{i<j}(r_i+r_j)$ is a nonzero
polynomial of degree at most
\[
 14+2+2+4=22
 \autoeqtag{eq:m9-selector-degree}
\]
in each root coordinate.  Consequently every binary $q\geq32$ admits five
distinct rational roots for which $Q$ is a genuine quadratic in $x$ with
$Q'\neq0$ and $Q\nmid R$.
\end{lemma}

\begin{proof}
Expanding $A_0A_3+A_1A_2$ with $A_j=B_j+xB_{j+1}$ gives
\eqref{eq:m9-Q-coefficients}; the two middle terms $B_1B_3$ cancel, and in
characteristic two $Q'=\mathrm dQ/\mathrm dx$ is the coefficient $q_1$.  Each
$A_j$ is linear in $x$, so $\Delta$, $N$ and $Q$ have degree at most two,
$N\Delta$ degree at most four, $(N\Delta)'$ degree at most two, and $R$
degree at most four.

Neither $q_2$ nor $q_1$ is the zero polynomial.  Written out,
\[
\begin{aligned}
 B_0&=z_4g_0+z_5g_1+z_6g_2+z_7g_3,&
 B_1&=z_3g_0+z_4g_1+z_5g_2+z_6g_3+z_7g_4,\\
 B_2&=z_2g_0+z_3g_1+z_4g_2+z_5g_3+z_6g_4+z_7,&
 B_3&=z_2g_1+z_3g_2+z_4g_3+z_5g_4+z_6,\\
 B_4&=z_2g_2+z_3g_3+z_4g_4+z_5 .
\end{aligned}
\]
In $q_1=B_0B_4+B_2^2$ the square $B_2^2$ contributes only square monomials.
The coefficients of $g_0^2$, $g_4^2$, $g_2^2$, $g_1^2$, $g_3^2$ and $1$ are
$z_2^2$, $z_6^2$, $z_4^2+z_2z_6$, $z_3^2$, $z_5^2+z_3z_7$ and $z_7^2$; read
in that order they force $z_2=z_6=z_4=z_3=z_5=z_7=0$.  In
$q_2=B_1B_4+B_2B_3$ the coefficients of $g_0g_1$, $g_1g_2$, $g_1g_4$ and
$g_3g_4$ are $z_2^2$, $z_3^2$, $z_4^2+z_3z_5+z_2z_6$ and $z_5^2+z_3z_7$; read
in that order they force $z_2=z_3=z_4=z_5=0$.  Both conclusions contradict
the hypothesis.

Pseudo-division by the degree-two polynomial $Q$ with leading coefficient
$q_2$ takes three steps, each replacing $R$ by $q_2R$ plus a multiple of $Q$;
carrying them out gives \eqref{eq:m9-pseudo-remainder}.  Over the fraction
field of $k[g_0,\ldots,g_4]$ the element $q_2$ is invertible, so $\rho=0$ would
say $Q\mid R$.  Lemma~\ref{lem:m9-nonsquare} rules that out: the double pole
along the component $\mathfrak q$ does not reduce, which is exactly the
nonvanishing of $R$ modulo $Q$.  Hence some $\rho_i\neq0$.

For the degrees, each $g_j$ is elementary symmetric in $r_1,\ldots,r_5$ and
therefore multi-affine, so each $B_j$ is multi-affine, each of
$q_2,q_1,q_0$ has
degree at most two in every root coordinate, each coefficient of $N\Delta$ at
most four, and each $R_i$ at most eight.  By
\eqref{eq:m9-pseudo-remainder} each $\rho_i$ has degree at most
$6+8=14$.  The Vandermonde $\prod_{i<j}(r_i+r_j)$ has degree four in each
root, which gives \eqref{eq:m9-selector-degree}.

Finally, a nonzero polynomial of degree at most $d$ in each of $n$ variables
separately does not vanish identically on $S^n$ for any set $S$ with
$|S|>d$: fix all but one variable at values where the polynomial is nonzero,
and use that a nonzero univariate polynomial of degree at most $d$ has at
most $d$ roots.  With $d=22$ and $S=\F_q$, every $q>22$, hence every binary
$q\geq32$, carries a point where the product above is nonzero.  There the
Vandermonde factor makes the five roots distinct, $q_2\neq0$ makes $Q$
quadratic in $x$, $q_1=Q'\neq0$, and $\rho\neq0$ gives $Q\nmid R$.
\end{proof}

\begin{theorem}[Empty first higher Lucas carrier]\label{thm:m9-shallow}
For every $m\geq4$,
\[
 \cM_9(\F_{2^m})\cap\{\textup{split-free syndromes}\}
 =\varnothing .
\]
Equivalently, every one of
\[
                  1+q+q^2+q^3+q^4+q^5
\]
projective carrier points has a split squarefree degree-eight member in its
Hankel kernel.
\end{theorem}

\begin{proof}
Proposition~\ref{prop:m9-invariant-block} handles $\PP(U)$.  Outside it,
projective transport normalizes the quotient in
\eqref{eq:lucas-m9-quotient} to $[e_4]$, leaving
the affine slice
\[
              (z_2,z_3,z_4,z_5,z_6,z_7)=(a,b,1,0,c,d)
\]
modulo its upper-Borel stabilizer.

In particular $z_4\neq0$, so Lemma~\ref{lem:m9-rank} applies:
\eqref{eq:binary-hankel-rank} has rank three or four, and in either case
Lemma~\ref{lem:m9-nonsquare} makes the final-pair cover
\eqref{eq:artin-schreier-cover} geometrically integral, with a component of
$Q=0$ along which the double pole does not reduce.  The rank split, the
derivative identity \eqref{eq:m9-square-obstruction} and its elimination are
manuscript calculations.  The public rank-two replay reconstructs the
complete \(q=64\) rational-twist mass and its five witnesses; it is not a
symbolic all-field proof.

Fix five roots and let $x$ be the sixth, as in
\eqref{eq:m9-five-root-split}.  Lemma~\ref{lem:m9-selector} then supplies,
for every binary $q\geq32$, five distinct rational roots at which $Q$ is
quadratic in $x$ with $Q'\neq0$ and $Q\nmid R$; its selector product has
individual root degree at most $22$, which is what the grid bound needs.
The reduced cover has at
most two simple poles, hence genus at most one.  The exact deletion table on the
normalized cover is
\[
\begin{array}{c|ccccc}
\text{condition}&x=r_i&Q=0&\Delta=0&
\text{final/old collision}&\text{final/}x\text{ collision}\\
\hline
\deg_x&5&2&2&10&4,
\end{array}
\]
so at most $48$ rational points are removed.  Since
$64+1-2\sqrt{64}=49$, every complement point is shallow for
$q\geq64$.

For $q=16,32$, the intrinsic Borel quotients contain respectively
$292$ and $1090$ orbits.  The compact certificates give one exact
eight-root witness for every orbit, and an independently written replay
reconstructs the quotient and both Hankel equations.  This closes the two
fields below the point-count threshold and proves the theorem.
\end{proof}

Projective-additive three-space divisors are sufficient but not necessary.
At $q=16$ their complete $\PGL_2$-closure has $510$ divisors, while
$101$ of the $292$ complement orbits contain none of them; ordinary
split witnesses still exist on all $101$.  Thus the endpoint criterion of
Wang--Wu--Hu and Proposition~\ref{prop:linearized} describe a proper section.
The exact uniform replacement is
\eqref{eq:final-pair-linear}--\eqref{eq:final-pair-quadratic}.

\begin{proposition}[Pointed R10 stage budgets]\label{prop:r10-pointed-budgets}
Let \(f\notin\cP_{10}\cup\cM^{\max}_{10,p}\), and retain \(i\) distinct
markers after \(i\) contractions, \(0\leq i\leq4\).  The unavailable
choices for the next marker form a zero-dimensional scheme of degree at
most \(19+i\), hence at most \(23\).  None of its defining factors vanishes identically off the
displayed carrier union.
\end{proposition}

\begin{proof}
At syndrome degree \(9-i\), the pullback of the lower persistent carrier
has degree at most three and the complete lower Lucas union has degree at
most two.  After its fixed factors are removed, the moving separable
linear series has ramification degree at most
\[
                         2(9-i)-4=14-2i.
\]
Each retained marker contributes one linear marker diagonal and one
fixed-marker gcd incidence cut out by quadratic minors, of degree at most
two.  Thus the full pointed degree is
\[
             3+2+(14-2i)+i+2i=19+i\leq23.
 \autoeqtag{eq:binary-deletion-budget}
\]
The persistent and Lucas pullbacks can vanish identically only on the
contained schemes classified by
Proposition~\ref{prop:recursive-component-selection}.  Identically zero
ramification is the inseparable case routed to the Lucas carrier.  A marker
diagonal is a proper point.  Finally, the evaluation-hyperplane argument in
the proof of Proposition~\ref{prop:r9-budgets} shows that an identically
vanishing fixed-marker minor forces the exact-linear-gcd package after
cancellation, whose contained case is persistent.  Hence every factor in
the bound \eqref{eq:binary-deletion-budget} is proper on the stated complement.
\end{proof}

\begin{proposition}[R10 fixed-depth escape]\label{prop:r10-escape}
For odd \(q\geq59\), and for even \(q\geq64\), every split-free
redundancy-ten syndrome belongs to
\(\cP_{10}\cup\cM^{\max}_{10,p}\).
\end{proposition}

\begin{proof}
Apply the finite-depth escape theorem through five markers.
Proposition~\ref{prop:r10-pointed-budgets} supplies all five successive
unavailable-marker schemes, including the retained-marker incidences omitted
by an unpointed iteration.  The terminal lower packages are the R8 package of
Propositions~\ref{prop:r8-bottom} and \ref{prop:r8-lp61}, followed by the
R9 package of Propositions~\ref{prop:r9-budgets} and
\ref{prop:r9-contained}; the R5--R7 terminal packages are those used in
their respective classification theorems.  At the top step, the generic
ordered-pair cover has genus at most one.  Its base diagonal and branch
delete \(12\) points, and each of the five forbidden markers deletes at
most \(6\), for the exact total
\[
                              12+5\cdot6=42.
\]
These are all terminal transverse, modular, and self-collision strata.
Since
\[
                              q+1-2\sqrt q>42
\]
first holds at a prime power \(q=59\), the odd assertion follows.
In characteristic two the ordered-Hessian base selector requires
\[
\min\{(9-4)(9+11)/2+1,\;9(9-4)\}=45
\]
parameters, and its reduced \((2,2)\) slice deletes
\(5+2+2+10+4=23\) points.  The first binary field satisfying both
conditions is \(64\).  The selector is obtained by pulling back one nonzero
equation from each of the common-quadratic and complementary Fano rulings:
their product has individual root degree eight, so Schwartz--Zippel or
nine-element block interpolation gives the displayed bound.  The normalized
\((2,2)\) ordered-Hessian slice has genus at most one, and its five deletion
degrees are exactly the displayed summands.  Proposition~\ref{prop:recursive-component-selection}
then identifies the sole contained locus in either case as the displayed
carrier union.
\end{proof}

\begin{proposition}[Characteristic-seven coherent lift at R10]
\label{prop:r10-char7-lift}
Let \(q=7^m\geq343\).  Every point of
\[
       \cM^{\max}_{10,7}\subseteq\PP\langle e_3,e_4,e_5,e_6\rangle
       \subset\PP(\Gamma^9E)
\]
has a split squarefree octic in its Hankel kernel.
\end{proposition}

\begin{proof}
Choose the marker at infinity.  In the affine-coordinate convention of
Section~\ref{sec:polar}, contraction deletes the last coordinate, so a
nonzero \(f\in\langle e_3,e_4,e_5,e_6\rangle\) contracts to the same
coefficient row
\[
 b=\iota_\infty f\in
 \PP\langle e_3,e_4,e_5,e_6\rangle
 \subset\PP\langle e_2,e_3,e_4,e_5,e_6\rangle,
\]
the characteristic-seven R9 binary-quartic carrier.  Since \(q>102\),
Proposition~\ref{prop:r9-components} selects four distinct finite base
roots.  The construction in Proposition~\ref{prop:r9-slice} then fixes a
fifth finite root as its moving parameter; because
\(q+1-2\sqrt q>32\), Proposition~\ref{prop:r9-deletion} selects a point
whose residual quadratic is separable and avoids all five fixed roots.
Its nonzero leading coefficient is part of the deleted determinant locus,
so both residual roots are finite.  The resulting split squarefree septic
\(g\in W_b\) therefore avoids the marked root at infinity.  Equation
\eqref{eq:lift} gives \(\lambda_\infty g\in W_f\), and marker avoidance
makes this octic squarefree.
\end{proof}

\begin{corollary}[Redundancy ten]\label{cor:r10}
For every prime power $q\geq59$, the projective deep-hole syndromes of
$\PRS(q-9)$ are exactly the persistent tangent and conjugate-secant
families.  Their cardinality is $q(q+1)^2/2$, and their orbit law is
$T/T^9$ modulo inversion and Frobenius, with tangent cocycle
$z\mapsto z+9u$.
\end{corollary}

\begin{proof}
For odd $q\geq59$, Proposition~\ref{prop:r10-escape} confines a
nonpersistent split-free point to the maximal Lucas carrier.  Its support
is empty except in characteristic seven; the
first field in this range is \(343\), where
Proposition~\ref{prop:r10-char7-lift} applies.
For even prime powers in the range one has $q\geq64$;
Proposition~\ref{prop:r10-escape} confines every split-free point to
$\cP_{10}\cup\cM_9$, and Theorem~\ref{thm:m9-shallow} eliminates the
second term.  Proposition~\ref{prop:upper-radius} holds throughout the stated range, so the
split-free equality promotes to deep holes.  The persistent count and orbit
law are the specialization $r=10$ of
Theorem~\ref{thm:recursive-carrier}.
\end{proof}
